\section*{Conclusion and Future Work}

This work presented a configurable VR sandbox for comparing several prism-effect implementations inside one shared baseline--exposure--post workflow. The prototype implemented translation, rotation, and skew perturbations; supported controller and embodied hand-tracking input; logged event and movement data; and provided an RShiny dashboard for condition, participant, quality, spatial, and trajectory analysis. In a 10-participant healthy evaluation, the system produced analyzable data across 80 completed modules. The clearest pattern was found in controller-based Line Bisection, where perturbation conditions were above participant-specific None baselines, but the broader dataset also showed non-zero no-effect drift and substantial participant variability.

The main conclusion is that the sandbox is useful as a framework-building and evaluation tool: it can place otherwise scattered VR prism-effect implementations into one comparable procedure and expose the procedural variables that shape interpretation. Future work should test a larger sample, add stricter posture and movement constraints, compare concurrent and terminal feedback directly, introduce stronger washout or reset procedures between modules, and improve embodied tracking robustness. SSQ should remain part of future evaluations, but should be paired with a separate fatigue or workload measure because raised-arm tiredness was more visible than simulator sickness in several sessions. Later rehabilitation-oriented studies can then test whether the same framework supports patient-facing procedures and repeated-session outcomes.
